% TOP_PROBLEM: {"id": 30006484, "title": "Open Problems on Spectral-Parameter Yang–Baxter Solutions", "category_id": 16, "category": {"id": 16, "name": "physics", "display_name": "Mathematical Physics", "description": "Problems at the intersection of mathematics and physics.", "slug": "physics", "order_index": 16, "created_at": "2026-07-31T15:26:25.671Z"}, "status": "open", "statusQualification": "For infinite-dimensional $V$, domains and the topology of invertibility must be specified. Different checked and unchecked $R$ conventions give different written equations, so the report's convention is retained."}
% FIRST_POSTED: 2026-10-11
% ENTRY_KIND: statement_only
% SUBMISSION: {"kind": "statement_only", "category": "top_problems", "problem_id": "30006484", "proposal_type": "established", "source": "Mini-Workshop: The Yang–Baxter Equation and Representations of Braid Groups (2025), Oberwolfach Reports, DOI 10.4171/owr/2025/49, pp. 2633--2672. Original report and mathematical context. Incorporated report material remains under the source CC BY-SA 4.0 terms unless otherwise noted. https://doi.org/10.4171/owr/2025/49", "status_evidence": "Source status retained from the supplied canonical registry. Dated literature assessment (2026-08-22): The listed structural, categorical, and spectral-parameter questions for Yang--Baxter solutions and Nichols algebras are posed as broad open directions; no complete classification or general theory resolving them was found. Primary report inspected for formulation; no comprehensive current-status literature review is claimed.", "status_checked_on": "2026-10-11", "duplicate_check": "Canonical numeric ID and public code checked against full registry; absent from baseline selected website registry. Source scope compared with adjacent source questions and existing related catalogue entries. No new identity allocated.", "attribution": "Original authors and contributors credited in Mini-Workshop: The Yang–Baxter Equation and Representations of Braid Groups (2026), Oberwolfach Reports, DOI 10.4171/owr/2025/49, pp. 2633--2672; UnsolvedMath curation credited separately. Definitions and English exposition prepared with Codex assistance; no independent human mathematical verification claimed."}
% SOURCE_STATEMENT: "Let $V$ be a possibly infinite-dimensional vector space and let $R:\\mathbb R\\to\\operatorname{GL}(V\\otimes V)$ satisfy the Yang–Baxter equation with spectral parameter. Can such families be studied by a unified algebraic or categorical theory, and how do analytic properties of $R(x)$ control the associated braided or quasi-braided categories?"
% Preparation and final self-audit: CONTRIBUTING.md prompts C and D; source/audit artifacts retained in the local physics_statements_2026_10_11 batch.
% Imported from: physics_statements.tex; UnsolvedMath ID 30006484
% !TeX program = lualatex
\documentclass[11pt]{article}
\usepackage[margin=25mm]{geometry}
\usepackage{fontspec}
\setmainfont{Latin Modern Roman}
\usepackage{amsmath,amssymb,mathtools}
\usepackage{unicode-math}
\setmathfont{Latin Modern Math}
\usepackage{xurl,hyperref}
\hypersetup{colorlinks=true,urlcolor=blue}
\setcounter{secnumdepth}{0}
\setlength{\parindent}{0pt}
\setlength{\parskip}{5pt}
\setlength{\emergencystretch}{3em}
\newcommand{\sref}[2]{\hyperlink{source:#1:#2}{[S#2]}}
\newcommand{\cataloguescope}[1]{\paragraph{Recorded target.}#1}
\newcommand{\R}{\mathbb R}
\newcommand{\C}{\mathbb C}
\newcommand{\N}{\mathbb N}
\newcommand{\Z}{\mathbb Z}
\newcommand{\Q}{\mathbb Q}
\newcommand{\F}{\mathbb F}
\newcommand{\eps}{\varepsilon}
\begin{document}
% UnsolvedMath ID 30006484; source catalogue code OWR-14299580-001
\setcounter{section}{30006483}
\section{Open Problems on Spectral-Parameter Yang–Baxter Solutions}\label{30006484}
{\small UnsolvedMath ID 30006484 · Mathematical Physics}
\subsection{Definitions and mathematical statement}

\paragraph{Shared notation.}$\mathbb N=\{1,2,\ldots\}$; the standard number fields have their usual meanings. All additional conventions are stated locally.


A spectral-parameter Yang--Baxter operator is an invertible map $R(x)$ on $V\otimes V$. Set $R_1(x)=R(x)\otimes I$ and $R_2(x)=I\otimes R(x)$. In the report's checked convention it satisfies, for every real $x,y$,
\[
R_1(x)R_2(x+y)R_1(y)=R_2(y)R_1(x+y)R_2(x).
\]
A braided category packages compatible exchanges; a quasi-braided variant allows the source's additional associativity data.

\paragraph{Catalogue wording retained for traceability.}
Let $V$ be a possibly infinite-dimensional vector space and let $R:\mathbb R\to\operatorname{GL}(V\otimes V)$ satisfy the Yang–Baxter equation with spectral parameter. Can such families be studied by a unified algebraic or categorical theory, and how do analytic properties of $R(x)$ control the associated braided or quasi-braided categories?

\paragraph{Scope and source qualification.}For infinite-dimensional $V$, domains and the topology of invertibility must be specified. Different checked and unchecked $R$ conventions give different written equations, so the report's convention is retained.
\subsection{Short English statement}
Develop a common theory of spectral-parameter Yang--Baxter families and relate their analytic behavior to braided or quasi-braided categories.
\subsection{Sources}
\begin{itemize}
\item[\hypertarget{source:30006484:1}{S1}] ULAM.AI and the UnsolvedMath Contributors, UnsolvedMath, record 30006484 (OWR-14299580-001). Catalogue curation under CC BY 4.0; definitions and English exposition added locally.
\url{https://huggingface.co/datasets/ulamai/UnsolvedMath}.
\item[\hypertarget{source:30006484:2}{S2}] Mini-Workshop: The Yang–Baxter Equation and Representations of Braid Groups (2025), Oberwolfach Reports, DOI 10.4171/owr/2025/49, pp. 2633--2672. Original report and mathematical context. Incorporated report material remains under the source CC BY-SA 4.0 terms unless otherwise noted.
\url{https://doi.org/10.4171/owr/2025/49}.
\end{itemize}
\label{end:30006484}
\end{document}
